\documentclass[supplement,review,onefignum,onetabnum]{siamart251216}

% Mathematical symbols. The SIAM class already loads amsmath.
\usepackage{amssymb}

% Tables.
\usepackage{booktabs}

% Oxford comma in multiple cleveref references.
\newcommand{\creflastconjunction}{, and~}

% Custom mathematical notation.
\newcommand{\fl}{\operatorname{fl}}
\newcommand{\R}{\mathbb{R}}

\emergencystretch=1.5em

\newcommand{\submissiondate}{August 03, 2026}

\title{Supplementary Materials: Backward Error of the Matrix Inverse in the Purcell--Egerv\'ary--Castillo Pivoting Transformation%
\thanks{Submitted to the editors \submissiondate.}}

% Running headers. The abbreviated title must not exceed 50 characters.
\headers{Supplementary Materials}{F. R. Villatoro}

\author{Francisco R. Villatoro%
\thanks{Escuela de Ingenier\'ias Industriales, Universidad de M\'alaga,
29071 M\'alaga, Spain
(\email{frvillatoro@uma.es}).}}

\begin{document}

\maketitle

These supplementary materials collect historical, technical, and methodological detail for the
main paper: the extended genealogy of the method (SM1), proof details omitted from the main
floating-point analysis (SM2), the design and provenance of the two campaigns (SM3), reliability
accounting (SM4), the empirical pivot law by family and variant (SM5), predictor comparison and
regression (SM6), per-family results of the large-scale campaign (SM7), the defect replay and its
masking decomposition (SM8), the solution backward-error reversal (SM9), the adversarial families
(SM10), high-precision validation of the decomposition (SM11), and the reproducibility archive and retained bulk outputs (SM12). Notation follows the
main paper. For a row-pivoted variant let $\Pi$ be the final processed-row permutation, set
$\bar A=\Pi A$ and $\bar b=\Pi b$, and let $T$ be the terminal tableau. The inverse represented for
the original matrix is $X=T\Pi$, with
\[
I-AX=\Pi^T(I-\bar A T)\Pi,
\qquad
(I-AX)X^{-1}=\Pi^T(I-\bar A T)T^{-1},
\qquad Xb=T\bar b.
\]
Thus the spectral, infinity, and Frobenius diagnostics used here are unchanged by working with the
row-permuted problem. As in the main paper, we suppress the bars in theoretical formulas and write
$\widehat V_{N+1}$ for the computed terminal tableau, restoring $\Pi$ explicitly in
implementation-facing identities when useful. With this convention, $u$ is the unit roundoff,
$\eta_{\mathrm{inv}}=\|(I-A\widehat V_{N+1})\widehat V_{N+1}^{-1}\|_2/\|A\|_2$ is the normwise
inverse backward error,
$r_R=\|I-A\widehat V_{N+1}\|_2/(\|A\|_2\|\widehat V_{N+1}\|_2)$ is the scaled right inverse
residual, $\eta_{\rm sol}=\|b-A\widehat x\|_2/(\|A\|_2\|\widehat x\|_2+\|b\|_2)$ is the normwise
solution backward error for $\widehat x=\fl(\widehat V_{N+1}b)$,
$D_N=I-A\widetilde V_{N+1}$ is the transformation defect, and $p^{\rm sc}_{\min}$ is the minimum
scaled pivot. The five pivoting variants \texttt{R0\_C0}, \texttt{R0\_C1}, \texttt{R0\_C2},
\texttt{R1\_C1}, and \texttt{R2\_C2} are defined in the main paper. We avoid hard-coded theorem
numbers below so that the supplement remains synchronized with revisions of the main manuscript.

\section{Extended genealogy of the method}
\label{sm:genealogy}

This section records the parts of the genealogy that the introduction of the main
paper compresses. Throughout, ``the transformation'' refers to the pivoting
transformation~(1.1) of the main paper.

\paragraph{Purcell (1953) and Egerv\'ary (1956--1960)}
Purcell~\cite{Purcell1953} homogenizes $Ax=b$ by an artificial variable, starts
from the Cartesian unit vectors, and produces at each step a set of vectors
satisfying the first $i$ equations, the multipliers being ratios of inner
products exactly as in the transformation; this is precisely the construction of
\cite[\S5.9]{CastilloCJPC2000}. Forsythe~\cite{Forsythe1954}, reviewing
Purcell's paper, attributed the idea also to Motzkin. Egerv\'ary showed
in~\cite{Egervary1957} that the vector method is a particular case of his rank
reduction procedure, and in the posthumous paper~\cite{Egervary1960} gave the
general matrix form
%
\begin{equation}\label{eq:sm-rank-reduction}
\widetilde A = A - A X \bigl(Y^T A X\bigr)^{-1} Y^T A ,
\end{equation}
%
together with two observations that are directly relevant to the main paper:
the elementary transformations are \emph{oblique projectors}, since
$A\bigl(I - X(Y^TAX)^{-1}Y^TA\bigr) = \bigl(I - AX(Y^TAX)^{-1}Y^T\bigr)A$, and
the initial matrix need not be the identity, which is the situation of
Example~3 of~\cite{CastilloCJPC2000}. A modern account of the rank reduction
theory is given by Gal\'antai~\cite{Galantai2009,Galantai2001}, and its
connection with Purcell's method is reviewed by
Heged\H{u}s~\cite{Hegedus2010}.

\paragraph{Hestenes (1958)}
Beyond the results quoted in the main paper, Theorem~6.6
of~\cite{Hestenes1958} shows that the deviation of the multiplier matrix from
the identity measures the departure from biorthogonality, which his \S7 turns
into an inexpensive running diagnostic; and his \S\S7--8 contrast the fragile
choice $V^{(0)}=I$, for which the pivots are ratios of leading principal minors
of $A$ itself, with the choice $V^{(0)}=U^*$, which yields a Gram matrix whose
leading principal minors are positive (Theorem~6.5).

\paragraph{The ABS class (1984--)}
Abaffy, Broyden, and Spedicato~\cite{AbaffyBS1984} placed the method inside a
general class of direct projection algorithms, introducing the term
\emph{implicit LU factors} and unifying a large family of finitely terminating
methods~\cite{AbaffySpedicato1989}. Within the ABS framework the numerical
behavior has been studied along two lines: sensitivity analysis in the model of
Broyden~\cite{Broyden1974,Broyden1985}, extended by
Gal\'antai~\cite{Galantai1991} and continued for two-step variants
in~\cite{AsadbeigiPBJ2017}, and forward error analysis in floating-point
arithmetic~\cite[Ch.~12]{AbaffySpedicato1989}, which yields growth functions
involving $\kappa(A)$ that the authors themselves describe as complex and
unduly pessimistic. Neither line produces a backward error bound. The ``twice
is enough'' criterion of Parlett and Kahan~\cite{Parlett1980}, in the modified
form due to Heged\H{u}s, has been transferred to the ABS classes by Abaffy and
Fodor~\cite{AbaffyFodor2015}, providing a practical, selective alternative to
unconditional reorthogonalization.

\paragraph{Benzi, Meyer, and T\r{u}ma (1993--1998)}
Benzi and Meyer~\cite{BenziMeyer1995}, following Benzi's
thesis~\cite{Benzi1993} and the independent work of T\r{u}ma~\cite{Tuma1993},
showed that, started from the identity, the process is mathematically
equivalent to ``right'' Gaussian elimination and produces Crout's form of the
LU decomposition. Benzi and T\r{u}ma~\cite{BenziTuma1998} extended the
construction to the nonsymmetric case by biconjugation, recovering the pivot
characterization $p_i = \Delta_i/\Delta_{i-1}$ already present
in~\cite{Hestenes1958}. The discrepancy between the computed $w_i^TAz_i$ and
$z_i^TA^Tw_i$ had been observed much earlier by Fox, Huskey, and
Wilkinson~\cite{FoxHW1948}, who attribute it to rounding and choose to retain
rather than discard the off-diagonal terms. The symmetric positive definite
case and the resulting AINV preconditioner are treated
in~\cite{BenziMeyerTuma1996}, its breakdown-free stabilization
in~\cite{BenziCullumTuma2000}.

\paragraph{Perturbation bounds}
Gal\'antai~\cite{Galantai2003,Galantai2005} derives exact perturbation
expressions for the LU, LDU, and Cholesky factorizations, applies them to the
full rank factorizations produced by Egerv\'ary's rank reduction procedure, and
concludes that the latter is stable whenever the LDU factorization of $Y^TAX$
is stable. These are the sharpest available results on the method; they
quantify the sensitivity of the computed factors to changes in the data, not
the existence of a nearby matrix of which the computed inverse is the exact
inverse.

\section{Proof details omitted from the floating-point analysis}
\label{sm:proofs}

\paragraph{Compact form of the local formation defect}
Use the permuted-column notation of the main paper,
$\overline V_j=\widehat V_jP_j$ and
$\overline v_k^{\,j}=\overline V_je_k$, and put
$q_k=|a_j|^T|\overline v_k^{\,j}|$. The componentwise estimates are
%
\[
|\varphi_{j,j}|
\le u+\gamma_n q_j|\hat\beta|,
\qquad
|\varphi_{j,k}|
\le\gamma_{n+1}q_k+\gamma_n q_j|\hat\alpha_k|.
\]
%
The $k$th column of $|\widehat B_j|$ is
$e_k+|\hat\alpha_k|e_j$ for $k\ne j$ and $|\hat\beta|e_j$ for $k=j$. The off-pivot
components are therefore bounded by
$\gamma_{n+1}|a_j|^T|\overline V_j|\,|\widehat B_j|$. For the pivot component no separate
relative-accuracy hypothesis is required. If
$\overline t_j^{\,j}$ is the selected computed pivot in permuted order, then
$|\overline t_j^{\,j}|\le(1+\gamma_n)q_j$ and
$|\hat\beta|\ge(1-u)/|\overline t_j^{\,j}|$, whence
%
\[
q_j|\hat\beta|\ge\frac{1-u}{1+\gamma_n}=(1-u)(1-nu).
\]
%
Consequently, if $(n+1)u<1$,
%
\[
u+\gamma_n q_j|\hat\beta|
\le
\left(\gamma_n+\frac{u}{(1-u)(1-nu)}\right)q_j|\hat\beta|
\le\gamma_{n+1}q_j|\hat\beta|,
\]
%
because
%
\[
\gamma_{n+1}-\gamma_n-
\frac{u}{(1-u)(1-nu)}
=
\frac{nu^2}{(1-u)(1-nu)(1-(n+1)u)}\ge0.
\]
%
Thus
$|\varphi_j^T|\le\gamma_{n+1}|a_j|^T|\widehat V_jP_j|\,|\widehat B_j|$.
Notice that the absolute value on $\hat\beta$ is essential because the pivot can be negative.

\paragraph{The cross term of the telescoped defect}
The second sum in the telescoped representation of $D_N$ vanishes when the transformations are
applied exactly and in general satisfies
%
\[
\left\|\sum_{j=1}^{N}e_j a_j^TE_j\widehat{\mathcal B}_{j:N}\right\|_2
\le
\gamma_2\max_j\|a_j\|_2
\sum_{j=2}^{N}\|\widehat{\mathcal B}_{j:N}\|_2
\sum_{l=1}^{j-1}
\bigl\|\,|\widehat V_lP_l|\,|\widehat B_l|\,\bigr\|_2
\|\widehat{\mathcal B}_{l+1:j-1}\|_2,
\]
%
by applying the accumulation bound to each partial error
$E_j=\sum_{l<j}F_l\widehat{\mathcal B}_{l+1:j-1}$. Thus the formation and application mechanisms
are algebraically distinct but coupled; they should not be described as independent sources.

\paragraph{Proof of the individual-solution bound}
Assume $b\ne0$. For a floating-point matrix--vector product there is a $\Delta V$ satisfying
$|\Delta V|\le\gamma_n|\widehat V_{N+1}|$ such that
$\widehat x=(\widehat V_{N+1}+\Delta V)b$. Hence
%
\[
b-A\widehat x=R_Nb-A\Delta Vb.
\]
%
Since the denominator of $\eta_{\rm sol}$ is at least $\|b\|_2$,
%
\[
\eta_{\rm sol}
\le
\frac{\|R_Nb\|_2}{\|b\|_2}
+
\gamma_n\|A\|_2\bigl\|\,|\widehat V_{N+1}|\,\bigr\|_2.
\]
%
Moreover $\Delta A=R_N\widehat V_{N+1}^{-1}$ implies
$R_N=\Delta A\widehat V_{N+1}$, and therefore
%
\[
\frac{\|R_Nb\|_2}{\|b\|_2}
\le
\eta_{\rm inv}\,\|A\|_2\,\|\widehat V_{N+1}\|_2.
\]
%
This makes explicit the conditioning factor in any uniform comparison between solution and inverse
backward error.

\paragraph{Full proof of the diagonal-dominance growth bound}
Let $A$ be strictly row-diagonally dominant with
$\|A\|_\infty=1$ and gap $\alpha>0$. For $i\le j$,
%
\[
|a_{ii}|-\sum_{k\le j,\,k\ne i}|a_{ik}|
\ge
|a_{ii}|-\sum_{k\ne i}|a_{ik}|
\ge\alpha,
\]
%
so every leading principal submatrix $A_j$ inherits the same gap. If $A_jx=b$ and
$|x_i|=\|x\|_\infty$, then
$|b_i|\ge\alpha\|x\|_\infty$, proving
$\|A_j^{-1}\|_\infty\le1/\alpha$.

One step of elimination replaces, for $i,k\ge2$,
$a_{ik}$ by $a'_{ik}=a_{ik}-(a_{i1}/a_{11})a_{1k}$. Using
$|a_{11}|-\sigma_1\ge\alpha$ gives
%
\[
|a'_{ii}|-\sum_{k\ge2,\,k\ne i}|a'_{ik}|
\ge
(|a_{ii}|-\sigma_i)
+
|a_{i1}|\left(1-\frac{\sigma_1}{|a_{11}|}\right)
\ge\alpha.
\]
%
Hence every Schur complement has the same positive gap. Its active pivot row therefore satisfies
$|t_l^{\,l}|\ge\alpha$ and, for every $k>l$,
$|t_l^{\,l}|-|t_k^{\,l}|\ge |t_l^{\,l}|-\sum_{m>l}|t_m^{\,l}|\ge\alpha$.
Thus both the first-nonzero and maximum-modulus column rules select the diagonal pivot and perform
no column interchange.

The explicit-tableau identity gives
%
\[
\|V^{j+1}\|_\infty
\le
\max\{1,\|A_j^{-1}\|_\infty(1+\|A_j'\|_\infty)\}
\le\frac{2}{\alpha},
\qquad
\|(V^{j+1})^{-1}\|_\infty\le1.
\]
%
Thus $\|t\|_1\le2/\alpha$, and the nontrivial row of $B_l$ has $1$-norm at most
$(1+2/\alpha)/\alpha\le3/\alpha^2$. Finally,
$\mathcal B_{l+1:N}=(V^{l+1})^{-1}V^{N+1}$ in exact arithmetic, so the terminal norm cancels the
denominator of the growth factor and
%
\[
\Gamma_N
\le
N\frac{2}{\alpha}\frac{3}{\alpha^2}
=
\frac{6n}{\alpha^3}.
\]

\paragraph{From the exact path to the first-order floating-point residual}
For fixed $A$ and $n$, the proof above gives the quantitative separations
$|t_l^{\,l}|\ge\alpha$ and
$|t_l^{\,l}|-|t_k^{\,l}|\ge\alpha$ for every admissible competitor $k>l$. Excluding underflow and
overflow, a finite-step induction along the exact path gives
$\widehat V_j=V^j+O(u)$ and computed dot products equal to their exact-path values plus $O(u)$.
For all sufficiently small $u$ these perturbations are smaller than the fixed positive pivot and
comparison margins, so the computed first-nonzero and maximum-modulus rules follow the exact path,
and $\widehat B_j=B_j+O(u)$ at every step. Hence the computed growth factor is
$\Gamma_N+O(u)$. Combining this with
$\gamma_{n+3}=(n+3)u+O(u^2)$ gives
%
\[
r_R^\infty
\le
\frac{6n(n+3)}{\alpha^3}u+O(u^2),
\qquad u\to0.
\]
%
The hidden constant is allowed to depend on the fixed matrix and on $n$; no uniform finite-$u$
path-preservation theorem is claimed. The ``first-nonzero'' rule here is the rule of the main paper,
which accepts the first computed pivot unequal to zero in the working format; a tolerance-based
variant would require a separate asymptotic specification of its threshold.

\section{Design of the two campaigns}
\label{sm:design}

Two campaigns are reported. As in the main analysis, the matrix $A$ denotes the array after
conversion to the working format; input-conversion error from an external real-valued matrix is not
folded into the reported arithmetic error. The \emph{mechanistic} campaign, described in the
remainder of this section, was designed to identify which quantity governs the backward error. The
\emph{large-scale} campaign was designed to test whether those conclusions survive dimension,
structure and massive sampling, and to measure the quantities the analysis defines but does not
bound. Its design is factorial and balanced by construction: fourteen matrix families, the
thirteen dimensions $n\in\{8,11,16,23,32,45,64,91,128,181,256,362,512\}$, both precisions, the
five variants of Table~1 of the main paper, and three refinement levels, with an identical
multiset of generating parameters at every $n$ so that no aggregate mixes populations. Stochastic
families receive $20\,000$ realizations per cell, of which a balanced subsample of exactly $2000$
is retained in every dimension; deterministic families, for which repeated realizations carry no
information, receive one matrix per (parameter, $n$) pair and are reported separately as
structured and extremal test cases, never as a sample. Seeds are the hash of
(family, $n$, replica) and are recorded. The unit roundoff is fixed throughout at $u_{32}=2^{-24}$
and $u_{64}=2^{-53}$ and stored with every record. Diagnostics of the inverse, which do not
depend on the right-hand side, are stored in a table keyed by (matrix, variant, precision);
diagnostics of the solution are stored separately, keyed additionally by right-hand side and
refinement level. Two families are constructed specifically for the diagonal-dominance theorem of the main paper,
at the reference level with $\|A\|_\infty=1$ and nominal dominance gap on the grid
$\alpha\in\{10^{-3},2\!\cdot\!10^{-3},5\!\cdot\!10^{-3},10^{-2},2\!\cdot\!10^{-2},
5\!\cdot\!10^{-2},0.1,0.2,0.3,0.5\}$: random sign patterns, and
\[
A=\frac{I-\beta P}{1+\beta},\qquad \beta=\frac{1-\alpha}{1+\alpha},
\]
where $P$ is zero-diagonal and row-stochastic, with positive off-diagonal entries. The latter formula gives row norm one and
gap $(1-\beta)/(1+\beta)=\alpha$ algebraically. The measured gap of the binary64 reference matrix
is stored as \texttt{alpha\_measured}; conversion to binary32 or binary64 working arrays is a
separate floating-point operation, so finite-precision theorem audits are reported as numerical
checks rather than exact identities. The code, generation manifests, canonical post-processing scripts, and the result summaries needed
to reproduce the paper will be released in a public GitHub repository accompanying the manuscript.
The release will record the exact commit hash and cryptographic digests of the canonical artifacts;
all reported statistics refer to this single recovered campaign and its explicitly identified
follow-up audits.

\paragraph{Independent DD regression verifier}
After synchronizing the residual constant with the $\gamma_{n+3}$ bound of the revised analysis, a
standalone binary64 verifier was run on four DD constructions, the four dimensions
$n\in\{8,16,32,64\}$, and six gaps $\alpha\in\{0.5,0.2,0.1,0.05,0.02,0.01\}$, for $96$ cases.
All regression checks passed. The largest observed ratios were $4.126249\times10^{-2}$ for
$\Gamma_N/(6n/\alpha^3)$ and $1.520482\times10^{-2}$ for
$r_R^\infty/(\gamma_{n+3}\Gamma_N)$, and no column interchange occurred. This is an independent
implementation regression check, not a numerical proof of the theorem.

The campaign comprises $18\,744$ matrix/right-hand-side configurations,
evaluated in binary32 and binary64 arithmetic, with zero, one, and two steps of
iterative refinement, and a residual accumulated in extended precision. All
computations were carried out with a single implementation in which the
elementary transformations, the column permutations, and the intermediate
tableaux are recorded, so that every diagnostic is derived from the same
computed quantities as the inverse itself. The reference methods are LU with
partial pivoting, Householder QR, the SVD, and Gauss--Jordan inversion with
partial and with complete pivoting; the Gauss--Jordan references are included
because they, like the pivoting transformation, form an explicit inverse.

\paragraph{Matrix families}
Seven families are used.

\begin{enumerate}
\item \texttt{random\_conditioned}: random matrices with prescribed singular
values, $A = U\Sigma W^T$ with $U$, $W$ random orthogonal and three singular
value profiles (geometric decay, one small singular value, and clustered
spectra) for $\kappa_2 \in \{1,10^2,10^4,\dots,10^{14}\}$,
$n\in\{8,16,32,64,128\}$, and twenty replicas per configuration. Since
$\|A\|_2=1$ by construction, $\kappa_2(A)=\|A^{-1}\|_2$ on this family.

\item \texttt{hilbert}: Hilbert matrices with $4\le n\le 20$.

\item \texttt{vandermonde}: Vandermonde matrices with equispaced, Chebyshev,
clustered, near-unit, and random nodes.

\item \texttt{kahan\_type}: Kahan-type triangular growth examples.

\item \texttt{adversarial\_rotation}: the orthogonal block-diagonal family of
the adversarial proposition of the main paper, with $2\times2$ orthogonal symmetric blocks
parametrized by $\varepsilon\in\{10^{-1},10^{-2},10^{-4},\dots,10^{-14}\}$.

\item \texttt{repeated\_bad\_pivots}: the geometric-decay variant of the same
family, described in section~\ref{sm:adversarial}.

\item \texttt{scaled\_adversarial}: row- and column-scaled versions of the
adversarial construction, designed to separate the absolute and scaled pivot
scores.
\end{enumerate}

\paragraph{Right-hand sides}
Right-hand sides are generated from a known solution in extended precision, so
that forward errors are measured against a reference accurate to well below the
working precision. Three types are used: random, alternating-sign, and the
right singular vector of $A$ associated with $\sigma_{\min}$, the last chosen
to probe the direction in which the residual operator is largest.

\paragraph{Completeness}
All $18\,744$ designed cases were reached, producing $552\,747$
records across methods and refinement levels with no failures of the case driver. Of the designed
cases, $99$ contain a nonfinite binary32 input matrix, generated parameters
whose entries overflow single precision, and are rejected before any solver
is called; the remaining $18\,645$ form the common subset on which all methods
are compared.

\section{Reliability of the measured quantities}
\label{sm:reliability}

\begin{table}
\centering
\caption{Paired ratio $r_R^{(32)}/r_R^{(64)}$ of the scaled right inverse residual over the
$44\,557$ matched configurations of the mechanistic campaign. The theoretical first-order scale
ratio is $u_{32}/u_{64}=5.37\times10^8$.}
\label{tab:sm-uscaling}
\begin{tabular}{lrrrr}
\hline
Variant & pairs & $q_{25}$ & median & $q_{75}$ \\
\hline
\texttt{R0\_C0} & $9\,093$ & $4.73\times10^{7}$ & $2.32\times10^{8}$ & $9.32\times10^{8}$ \\
\texttt{R0\_C1} & $8\,907$ & $2.52\times10^{8}$ & $4.48\times10^{8}$ & $5.90\times10^{8}$ \\
\texttt{R0\_C2} & $8\,888$ & $2.49\times10^{8}$ & $4.44\times10^{8}$ & $5.97\times10^{8}$ \\
\texttt{R1\_C1} & $8\,852$ & $2.65\times10^{8}$ & $4.51\times10^{8}$ & $5.75\times10^{8}$ \\
\texttt{R2\_C2} & $8\,817$ & $2.72\times10^{8}$ & $4.53\times10^{8}$ & $5.89\times10^{8}$ \\
\hline
\end{tabular}
\end{table}

The exact normwise inverse backward error requires a solve with the computed tableau. The campaign
does not treat an arbitrary solve through a numerically unresolved $\widehat V_{N+1}$ as a
measurement. Instead it applies the conservative screen described in the main paper. With $W$ the
separately maintained inverse of the accumulated product of the computed transformations, both
$W$ and $\widehat V_{N+1}$ are cast to binary64 and the implementation forms
%
\[
\delta_F
=
\|I-\fl_{64}(W\widehat V_{N+1})\|_F
+
\gamma_n^{(64)}\|W\|_F\|\widehat V_{N+1}\|_F.
\]
%
The second term in $\delta_F$ is the standard multiplication-roundoff allowance used by the
code. If the two displayed terms were evaluated exactly, the triangle inequality would give
$\|I-W\widehat V_{N+1}\|_2\le\delta_F$. Thus, for $\delta_F<1$,
$\widehat V_{N+1}^{-1}=(I-H)^{-1}W$ with $H=I-W\widehat V_{N+1}$, and hence
%
\[
\|\widehat V_{N+1}^{-1}\|_2
\le\frac{\|W\|_2}{1-\delta_F},
\qquad
\kappa_2(\widehat V_{N+1})
\le
\frac{\|\widehat V_{N+1}\|_F\|W\|_F}{1-\delta_F}.
\]
%
This motivates the implemented quantity
%
\[
\kappa^{\rm scr}
=
\frac{\|\widehat V_{N+1}\|_F\|W\|_F}{1-\delta_F}
\]
%
and the retention rule $u\kappa^{\rm scr}<10^{-2}$. The screen itself is evaluated in binary64
without directed rounding, so we refer to retained records as \emph{reliability-screened}, not as
formally certified evaluations of $\kappa_2(\widehat V_{N+1})$. This terminology also avoids
confusing the operational screen with the direct condition number of the input matrix used later to
define the well-resolved regime for solution comparisons.

Of the $91\,466$ successful unrefined Castillo records in the mechanistic campaign, $52\,184$
($57.1\%$) pass this screen. The proportion varies by less than $0.3$ percentage points across the
five variants. Of these, $792$ have $\eta_{\mathrm{inv}}=0$ and are absent from logarithmic
displays; all belong to the two adversarial families where the pivoted variants invert the matrix
exactly in the recorded arithmetic. The remaining $51\,392$ positive values are the population
used in the predictor comparison below. The scaled right residual $r_R$ requires no solve and is
reported for every successful record.

The screen is intentionally distinct from the condition
$\kappa_2(A)u<10^{-2}$ used to classify a linear system as well resolved when comparing
$\eta_{\rm sol}$. The former concerns the numerical evaluation of a diagnostic of the stored
inverse; the latter concerns whether the original linear system is resolvable in the working
precision.

\section{The empirical pivot law by family and variant}
\label{sm:law}

The numerical-evidence section of the main paper reports that the empirical inequality
$\eta_{\mathrm{inv}}\le u\,n/p^{\rm sc}_{\min}$ holds for $99.0\%$ of the
reliability-screened records. Table~\ref{tab:sm-law-by-family} gives the breakdown by matrix
family and by pivoting variant, in terms of the ratio
$\eta_{\mathrm{inv}}\,p^{\rm sc}_{\min}/(u\,n)$, a value of at most $1$ meaning
that the inequality holds. The failures are concentrated in the Vandermonde
family and, across variants, in the rules with weaker pivoting; the ordering by
variant reproduces that of every other inverse diagnostic in the main paper.

\begin{table}
\centering
\caption{Behavior of the ratio $\eta_{\mathrm{inv}}\,p^{\rm sc}_{\min}/(u\,n)$ by matrix
family and by variant. A value at most $1$ means that the empirical law of the predictor subsection of the main paper holds.}
\label{tab:sm-law-by-family}
\begin{tabular}{lrrrr}
\hline
Family & records & holds & $q_{99}$ & max \\
\hline
\texttt{adversarial\_rotation} & $1\,200$ & $100.0\%$ & $2.50\times10^{-1}$ & $2.80\times10^{-1}$ \\
\texttt{repeated\_bad\_pivots} & $1\,350$ & $100.0\%$ & $2.50\times10^{-1}$ & $2.50\times10^{-1}$ \\
\texttt{scaled\_adversarial}   & $3\,878$ & $100.0\%$ & $1.49\times10^{-1}$ & $2.50\times10^{-1}$ \\
\texttt{hilbert}               & $75$     & $100.0\%$ & $2.03\times10^{-1}$ & $2.03\times10^{-1}$ \\
\texttt{kahan\_type}           & $15$     & $100.0\%$ & $8.92\times10^{-1}$ & $8.92\times10^{-1}$ \\
\texttt{random\_conditioned}   & $45\,130$& $99.1\%$  & $9.33\times10^{-1}$ & $1.22\times10^{2}$ \\
\texttt{vandermonde}           & $536$    & $79.5\%$  & $1.32\times10^{3}$ & $3.77\times10^{3}$ \\
\hline
\texttt{R2\_C2} & --- & $100.0\%$ & $3.10\times10^{-1}$ & $8.92\times10^{-1}$ \\
\texttt{R1\_C1} & --- & $99.97\%$ & $2.89\times10^{-1}$ & $2.22$ \\
\texttt{R0\_C2} & --- & $98.5\%$  & $1.46$ & $9.48\times10^{2}$ \\
\texttt{R0\_C1} & --- & $98.4\%$  & $1.47$ & $1.81\times10^{3}$ \\
\texttt{R0\_C0} & --- & $98.1\%$  & $2.02$ & $3.77\times10^{3}$ \\
\hline
\end{tabular}
\end{table}

\section{Predictor comparison, regression, and limitations}
\label{sm:regression}

The main paper distinguishes the best ex-post diagnostic from the best local trajectory indicator.
Table~\ref{tab:sm-predictors} compares candidate predictors through the spread of
$\eta_{\mathrm{inv}}/\text{predictor}$ over the $51\,392$ reliability-screened records with
$\eta_{\mathrm{inv}}>0$; smaller spread means closer tracking.

\begin{table}
\centering
\caption{Ratio $\eta_{\mathrm{inv}}/\text{predictor}$ over the $51\,392$ screened records with
$\eta_{\mathrm{inv}}>0$. The last column is $\log_{10}(q_{99}/q_{01})$; smaller is tighter.}
\label{tab:sm-predictors}
\begin{tabular}{lrrrr}
\hline
Predictor & $q_{50}$ & $q_{90}$ & $q_{99}$ & spread \\
\hline
$\|D_N\|_2/\|A\|_2$ alone     & $3.40\times10^{-1}$ & $1.17$ & $3.26\times10^{3}$ & $6.9$ \\
$u\,n/p^{\rm sc}_{\min}$      & $5.19\times10^{-2}$ & $1.97\times10^{-1}$ & $1.02$ & $8.3$ \\
$u\,n\,\rho_\mu$              & $6.72\times10^{-2}$ & $3.48\times10^{-1}$ & $4.76\times10^{4}$ & $12.5$ \\
$u\,n\,\kappa_2(A)$           & $2.32\times10^{-2}$ & $4.24\times10^{-1}$ & $6.38\times10^{1}$ & $17.2$ \\
$\Lambda_N^{\rm corr}$        & $4.26\times10^{-6}$ & $9.45\times10^{-3}$ & $4.97\times10^{-2}$ & $19.8$ \\
\hline
\end{tabular}
\end{table}

Under this metric $D_N$ is the tightest ex-post diagnostic, not
$un/p^{\rm sc}_{\min}$. The role of the latter is different: it is local, inexpensive, and tied
directly to the formation bound in the main paper. Its empirical inequality for
$\eta_{\mathrm{inv}}$ holds in $99.0\%$ of the screened records, but the failure rate is
class-dependent and reaches $20.5\%$ on Vandermonde matrices. The main paper therefore does not
promote it to a theorem or call it a complete predictor.

To quantify the relative association of local and growth diagnostics with the observed error we
fit
%
\begin{equation}\label{eq:sm-regression}
\begin{aligned}
\log_{10}\eta_{\mathrm{inv}}=c_0
&+c_1\log u+c_2\log n+c_3\log\frac{1}{p^{\rm sc}_{\min}}
+c_4\log\max_j\|\widehat V_j\| \\
&+c_5\log\max_j\|\widehat{\mathcal B}_{j+1:N}\|
+c_6\log\kappa_2(\widehat V_{N+1}).
\end{aligned}
\end{equation}
%
The in-sample fit gives $R^2=0.925$, with $c_1=0.975$ and $c_3=0.940$. Nested fits give
$R^2=0.342$ for $(u,n)$ alone and $0.890$ after adding $p^{\rm sc}_{\min}$; the three growth
quantities together add $0.035$ in sample. These coefficients are descriptive, not structural.
Only two values of $u$ occur, the fitted coefficient of $\log n$ is negative, and the regressors
are strongly correlated.

Most importantly, the model does not generalize uniformly across matrix classes. Under
leave-one-family-out validation the out-of-sample $R^2$ is $0.969$ for Hilbert, $0.868$ for the
prescribed-singular-value family, and $0.829$ for Vandermonde, but $0.057$ for the single-pivot
orthogonal family, $-0.405$ for repeated bad pivots, and $-0.802$ for the scaled adversarial family.
The reduced model in $(u,n,p^{\rm sc}_{\min})$ can fare worse still, reaching $R^2=0.036$ on the
prescribed-singular-value family. Growth variables therefore contribute little to the in-sample
increment after the pivot is known but are not dispensable as class-dependent information. The
regression should be read as a map of the tested predictor space, not as a universal law.

A historical analysis field named \texttt{rho\_inv} is deliberately not interpreted here as
$\eta_{\mathrm{inv}}/(u\kappa_2(\widehat V_{N+1}))$. In the audit that field was formed with the
conservative conditioning-screen denominator and is therefore only a conservative proxy for that
ratio; no certified ordering relative to the exact ratio is claimed. The revised main paper consequently makes no claim that the explicit-inverse inequality
$\eta_{\mathrm{inv}}\le r_R\kappa_2(\widehat V_{N+1})$ is observed to be close to equality.

\section{The large-scale campaign: per-family results}
\label{sm:large}

The large-scale campaign covers fourteen matrix families and the thirteen dimensions
$n\in\{8,11,16,23,32,45,64,91,128,181,256,362,512\}$ in both precisions and for the five
variants of Table~1 of the main paper. After the wall-clock limit of the array job it
contains $14\,686\,763$ complete matrices, $1.468\times10^{8}$ inverse records and
$1.322\times10^{9}$ solution records, with a core of $20\,000$ realizations per cell complete for
$n\le181$ and a balanced extension of exactly $2000$ realizations per stochastic cell in all
thirteen dimensions. Cells of the balanced extension are the unit of every aggregate reported in
the main paper.

\paragraph{Breakdown}
The overall breakdown rate is $327\,222/1.468\times10^{8}=0.2228\%$, distributed as $85.6\%$
\texttt{no\_pivot}, $13.3\%$ nonfinite tableau, $0.94\%$ nonfinite $B_l$ and $0.18\%$ nonfinite
binary32 input. It is concentrated by precision and by family rather than by method: in binary64
every variant breaks down on fewer than $8\times10^{-6}$ of records, and $83.6\%$ of all
breakdowns come from the prescribed-singular-value family at $\kappa_2\ge10^{8}$, where
$\kappa_2u>1$ in binary32 and the problem is unresolved in the working precision. Vandermonde
contributes $14.4\%$ and the diagonally dominant and scaled adversarial families essentially
none. Breakdown rates should therefore be read together with the conditional residual, never
separately; the Vandermonde comparison of the pivot-rule comparison of the main paper is the clearest instance.

\paragraph{Ordering by family}
The median across cells of $q_{95}(r_R/u)$ at $n=512$ is, in binary32 and binary64,
$1.35$ and $1.48$ for \texttt{R1\_C1}, $1.48$ and $1.67$ for \texttt{R2\_C2}, $1.96$ and $2.10$
for \texttt{R0\_C1}, $2.20$ and $2.47$ for \texttt{R0\_C2}, and $9.94$ and $15.35$ for
\texttt{R0\_C0}. The pooled ordering, however, conceals a systematic dependence on the matrix
class. On the prescribed-singular-value family row pivoting helps, \texttt{R1\_C1} beating
\texttt{R0\_C1} on $70.7\%$ and $74.0\%$ of paired matrices in the two precisions. On
\texttt{row\_diagonal\_dominant\_random} it hurts, \texttt{R0\_C1} winning on $85.3\%$ and
$76.2\%$; on \texttt{row\_diagonal\_dominant\_stochastic} the figures are $88.9\%$ and $85.3\%$.
On Vandermonde at $n=512$ in binary32 the breakdown rates are $65.6\%$, $0\%$, $0.1\%$, $72.7\%$
and $84.6\%$ for the five variants in the order of Table~1, with conditional
$q_{95}(r_R/u)$ equal to $1.52\times10^{5}$, $0.793$, $0.288$, $0.488$ and $0.296$; restricted
to matrices on which both succeed, \texttt{R0\_C2} has the smaller residual than \texttt{R2\_C2}
on $59\%$ of cases in binary32 and $58\%$ in binary64, with median ratio $0.77$.

\paragraph{Effective slopes}
Power-law fits $q_{95}(r_R/u)\sim Cn^{s}$ are computed per cell over the thirteen dimensions of
the balanced extension, with $95\%$ intervals from $2000$ cluster-bootstrap replications that
resample matrices within each $n$ and never resample the dimensions. Median interval widths are
$0.011$--$0.014$ for the pivoted variants. Almost all intervals exclude zero, but not with a
common sign: for \texttt{R0\_C1} in binary32, $54$ of $64$ cells have an interval entirely in
$s>0$ and $10$ entirely in $s<0$, and the pattern is similar for the other variants, so the
heterogeneity between cells is signal and not bootstrap noise. Comparing the power model with a
quadratic in $\log n$ by AIC, the quadratic is preferred in $47$ to $60$ cells out of $64$
depending on variant and precision, with $\Delta\mathrm{AIC}<-10$ in about three quarters of
those; among the cells where it wins, the curvature is convex in the majority for all four
pivoted variants. The fitted $s$ must therefore be reported as an effective slope over
$8\le n\le512$.

\paragraph{Sample size and wall-clock truncation}
Comparing quantiles from all $20\,000$ realizations with those from the balanced $2000$ subset in
the complete dimensions, the median relative difference is $0.57\%$ at $q_{50}$, $0.86\%$ at
$q_{90}$, $1.06\%$ at $q_{95}$, and $1.95\%$ at $q_{99}$; excluding the scaled adversarial family,
where relative differences near zero are not meaningful, the $95$th percentile of the difference
at $q_{95}$ is $5.4\%$. This is a sample-size audit: it supports $2000$ realizations for estimating
quantiles through $q_{95}$ on complete cells and cautions against $q_{99}$.

It is not a missingness audit. The array wall-clock limit truncated the largest dimensions through
the deterministic task/manifest assignment, so the recovered $2000$ records at $n>181$ are not
asserted to be a random subsample of the designed $20\,000$. We therefore distinguish the complete
core $n\le181$ from the balanced recovered extension through $n=512$. The latter is used for
finite-range dimensional trends and explicitly paired comparisons, without treating the missingness as random and with no
claim that a small complete-cell subsampling error proves absence of truncation bias.

\section{The defect replay and masking decomposition}
\label{sm:replay}

The replay recomputes $\|D_N\|_2$, $\|R_N\|_2$, and $\|\Pi AE_{N+1}\|_2$ for the balanced sample
of the two diagonally dominant ensembles: $520\,000$ matrices under \texttt{R0\_C0},
\texttt{R0\_C1}, \texttt{R1\_C1}, and \texttt{R2\_C2} in both precisions, giving
$4\,160\,000$ method/precision records and $1\,040\,000$ complete
\texttt{R1\_C1}/\texttt{R0\_C1} pairs. The exactly applied tableau is accumulated in extended
precision from the rounded elementary transformations. The matrix identity
$R_N=D_N-\Pi AE_{N+1}$ passes the absolute long-double first-order acceptance gate implemented by
the canonical replay code on every replayed record. This gate is operational rather than an
interval certificate: it uses no directed rounding and is not invoked as a theorem assertion. The
earlier auxiliary diagnostic normalized by $\|R_N\|$ is not an acceptance criterion because it
produces false rejections in binary64 when $\|R_N\|$ is a cancellation residual near the roundoff
floor; that relative diagnostic can reach approximately $1.46\times10^{-3}$. The scalar
factorization identities below, which are different diagnostics, close to approximately
$6\times10^{-16}$. The exact acceptance threshold and source hashes are fixed by the canonical
replay script distributed with the computational archive.

\paragraph{Control}
On all $520$ blocks, \texttt{R0\_C0} and \texttt{R0\_C1} reproduce the same pivot path and yield
$D_N$, $R_N$, and $\Pi AE_{N+1}$ equal bit for bit, with zero path mismatches. This is the numerical
counterpart of the exact diagonal-dominance result in the main paper and a stringent internal
control of the replay.

\paragraph{Exact scalar factorization}
The ratios below are defined only when their displayed denominators are positive.  The exact
pairwise identities therefore apply on the common domain where all displayed factors are defined.
A pair with a zero denominator is omitted from any ratio that requires it. On this natural domain,
for each \texttt{R1\_C1}/\texttt{R0\_C1} pair define
%
\[
d=\frac{\|D_1\|_2}{\|D_0\|_2},
\qquad
a=\frac{\|AE_1\|_2}{\|AE_0\|_2},
\qquad
H_i=(\|D_i\|_2^2+\|AE_i\|_2^2)^{1/2},
\qquad
q=\frac{Q_1}{Q_0},
\]
%
where $Q_i=\|R_i\|_2/H_i$, and let
%
\[
w=\frac{\|D_0\|_2^2}{\|D_0\|_2^2+\|AE_0\|_2^2}.
\]
%
Then, with $h=H_1/H_0$ and $r=\|R_1\|_2/\|R_0\|_2$,
%
\[
r=hq,
\qquad
h^2=wd^2+(1-w)a^2.
\]
%
These are exact scalar identities built from spectral norms. In particular, $Q$ is not a matrix
angle and $H$ is not the norm of an orthogonal decomposition.

To isolate the baseline-weight effect from the actual change in the application term, introduce
the counterfactual scale
%
\[
h_{a=1}=\sqrt{wd^2+1-w},
\qquad
M_{\rm weight}=\frac{h_{a=1}}{d},
\qquad
M_{AE}=\frac{h}{h_{a=1}}.
\]
%
Then every pair satisfies the exact multiplicative factorization
%
\begin{equation}\label{eq:sm-compensation}
r=d\,M_{\rm weight}\,M_{AE}\,q.
\end{equation}
%
At the pair level,~\eqref{eq:sm-compensation} is exact.  A product of geometric means is an exact
group identity only when all factors are evaluated on the same common positive-domain subset.
Accordingly, the four-factor products displayed below are used as descriptive summaries of the
pairwise decomposition, not as an additional theorem; medians are never multiplied.

The geometric-mean factors for B1 are
%
\[
2.572\times0.474\times0.966\times0.927=1.092
\quad\text{(binary32)},
\]
%
and
%
\[
2.489\times0.492\times0.966\times0.929=1.098
\quad\text{(binary64)}.
\]
%
For B2 they are
%
\[
2.156\times0.658\times0.998\times0.982=1.391,
\qquad
2.127\times0.677\times0.999\times0.983=1.414.
\]
%
Thus the dominant attenuation of the formation-defect amplification is the baseline-weight factor
$M_{\rm weight}$; the application-term factor is essentially unity, and the $Q$ factor gives only a
secondary reduction.

\paragraph{Weight dependence and deciles}
The median baseline weight $w$ is $0.086$ and $0.109$ on B1 in binary32 and binary64, and $0.367$
and $0.414$ on B2. The correlation of $\log d$ with $\log w$ is $-0.77$ and $-0.83$ on B1 and
$-0.50$ and $-0.52$ on B2: matrices whose formation defect is amplified most strongly tend to be
those on which the defect carried less baseline quadratic weight. This is an association, not a
causal coefficient.

Deciles make the same point without relying on correlation. On B2 in binary64 at $n=512$, the
lowest-$w$ decile has $w=0.0048$, $d=6.14$, and $r=0.96$, whereas the highest decile has
$w=0.586$, the smaller $d=5.09$, and $r=3.97$. On B1 in binary64 at $n=512$, the corresponding
values are $w=0.0047$, $d=6.39$, $r=0.93$ and $w=0.0255$, $d=6.40$, $r=1.36$. Across dimensions
on B1, the defect ratio rises from $1.08$ at $n=8$ to $7.03$ at $n=512$ while the median baseline
weight falls from $0.443$ to $0.013$. The final residual therefore masks most strongly the
formation-defect deterioration in the large-$n$ regime where that deterioration is largest.

\section{Mechanism of the solution backward-error reversal}
\label{sm:reversal}

Taken over all successful cases, the $99$th percentile of $\eta_{\rm sol}$
ranks the first-nonzero variant \texttt{R0\_C0} as the best of the five,
in contradiction to every diagnostic of the computed inverse; the reversal
survives restriction to the configurations solved by all five variants, and
disappears on restriction to the well resolved regime $\kappa_2(A)\,u<10^{-2}$
(Table~\ref{tab:sm-etasol}). Two distinct mechanisms operate in the unresolved
regime, and it is worth separating them because only the first is usually
anticipated.

\begin{table}
\centering
\caption{Solution backward error over configurations solved by all five
variants, with zero refinement. Left: all such configurations. Right: those in
the well resolved regime $\kappa_2(A)\,u<10^{-2}$; this solution-reliability
criterion is distinct from the inverse-reliability screen of the main paper, which is based on the
separately maintained inverse of the transformation product.}
\label{tab:sm-etasol}
\begin{tabular}{lrrrr}
\hline
& \multicolumn{2}{c}{all, $n=17\,283$} & \multicolumn{2}{c}{well resolved, $n=10\,045$} \\
\cmidrule(lr){2-3}\cmidrule(lr){4-5}
Variant & median & $q_{99}$ & median & $q_{99}$ \\
\hline
\texttt{R0\_C0} & $2.80\times10^{-9}$ & $5.90\times10^{-2}$ & $1.84\times10^{-11}$ & $5.07\times10^{-4}$ \\
\texttt{R0\_C1} & $1.23\times10^{-9}$ & $2.90\times10^{-1}$ & $4.91\times10^{-12}$ & $1.16\times10^{-4}$ \\
\texttt{R0\_C2} & $1.20\times10^{-9}$ & $2.80\times10^{-1}$ & $4.99\times10^{-12}$ & $1.23\times10^{-4}$ \\
\texttt{R1\_C1} & $1.24\times10^{-9}$ & $2.97\times10^{-1}$ & $4.97\times10^{-12}$ & $1.16\times10^{-4}$ \\
\texttt{R2\_C2} & $1.21\times10^{-9}$ & $2.86\times10^{-1}$ & $4.89\times10^{-12}$ & $1.10\times10^{-4}$ \\
\hline
\end{tabular}
\end{table}

\paragraph{Denominator inflation}
For records with $\eta_{\rm sol}>0$, the denominator of $\eta_{\rm sol}$ can be recovered without
cancellation as $\|A\|_2\|\widehat x\|_2+\|b\|_2 =
\|b-A\widehat x\|_2/\eta_{\rm sol}$, both quantities being recorded.  Zero-residual records
are excluded from this ratio reconstruction. Comparing \texttt{R0\_C0} with \texttt{R1\_C1} over
the $17\,679$ paired configurations, the median denominator ratio is $1.000$
in the well resolved regime, with $q_{90}=1.10$. When the problem is resolved,
the two variants produce solutions of essentially the same norm and
$\eta_{\rm sol}$ compares like with like. In the unresolved regime the median
ratio is $1.001$ but $q_{90}$ rises to $681$, and among the $3\,684$
configurations in which \texttt{R0\_C0} attains the smaller $\eta_{\rm sol}$
the median denominator ratio is $1.17$ with $70.1\%$ exceeding unity.
Denominator inflation is therefore real, and it accounts for part of the
reversal, since a less accurate inverse can return a smaller backward error simply by
producing a larger $\|\widehat x\|$.

\paragraph{Accidental cancellation in the residual}
Inflation is not, however, the whole explanation. In the extreme tail, the
$177$ paired configurations in the top percentile of $\eta_{\rm sol}$ for
\texttt{R1\_C1}, the median denominator ratio is $0.83$, and only $44\%$ of
cases have an inflated denominator. What distinguishes these cases is the
residual itself, which is a median factor $26$ smaller for \texttt{R0\_C0}.
Both computed solutions are meaningless there, the forward errors having medians
$2.50$ and $3.02$, and in $65.5\%$ of the cases both exceed unity, so neither
solution has a single correct digit. In that situation the residual is governed
by accidental cancellation rather than by the quality of the inverse, and a
smaller residual carries no information about the method. That \texttt{R0\_C0}
produces the worse inverse in $94.9\%$ of these same cases, with median ratio of
scaled inverse residuals $9.2$, makes the point directly.

\paragraph{Characterization of the unresolved regime}
Of the cases with $\eta_{\rm sol}>10^{-3}$, ninety percent are binary32, the
median condition number is $2.3\times10^{8}$ so that $\kappa_2(A)\,u>1$, and
the median forward error is $1.27$. On the well resolved subset
$\kappa_2(A)\,u<10^{-2}$, \texttt{R0\_C0} is the worst variant at both the
median and the $99$th percentile of $\eta_{\rm sol}$, by factors of $3.8$ and
$4.4$ respectively; its forward error is a median factor $17.8$ larger, and it
produces the more accurate solution in only $5.7\%$ of paired comparisons.
All cross-method comparisons of $\eta_{\rm sol}$ in the main paper are
therefore restricted to $\kappa_2(A)\,u<10^{-2}$, a threshold fixed before the
campaign was run.

\section{The adversarial families}
\label{sm:adversarial}

The adversarial proposition of the main paper proves the behavior of the pivot rules on the exactly conditioned
orthogonal family. Table~\ref{tab:sm-adversarial} records the floating-point confirmation over
eight decades of $\varepsilon$: the tableau growth of \texttt{R0\_C0} agrees with
$\sqrt{1+c}/\varepsilon$ to all displayed digits, while the no-row-pivoting magnitude rules agree with the
exact value $(1-\varepsilon)^{-1/2}$ and hence tend to one as $\varepsilon\to0$.  The underlying
real-valued matrices are exactly orthogonal; their binary64 representations satisfy
$\kappa_2(A)=1+O(u)$.

\begin{table}
\centering
\caption{Orthogonal family of the adversarial proposition of the main paper, binary64, $n=8$. Medians over the
recorded cases.  The underlying real-valued matrices have $\kappa_2(A)=1$ exactly; their binary64
representations have $\kappa_2(A)=1+O(u)$. The tableau growth of \texttt{R0\_C0}
agrees with $\sqrt{1+c}/\varepsilon=(\sqrt2/\varepsilon)(1+O(\varepsilon^2))$, while the displayed
\texttt{R0\_C2} growth agrees with $(1-\varepsilon)^{-1/2}$ (for example $1.054\ldots$ at
$\varepsilon=10^{-1}$ and $1.005\ldots$ at $10^{-2}$).}
\label{tab:sm-adversarial}
\begin{tabular}{lrrrr}
\hline
& \multicolumn{2}{c}{\texttt{R0\_C0}} & \multicolumn{2}{c}{\texttt{R0\_C2}} \\
\cmidrule(lr){2-3}\cmidrule(lr){4-5}
$\varepsilon$ & $\max_j\|\widehat V_j\|_2$ & $r_R$ & $\max_j\|\widehat V_j\|_2$ & $r_R$ \\
\hline
$10^{-1}$  & $1.41\times10^{1}$  & $3.61\times10^{-16}$ & $1.05$ & $0$ \\
$10^{-2}$  & $1.41\times10^{2}$  & $5.12\times10^{-15}$ & $1.01$ & $2.22\times10^{-16}$ \\
$10^{-4}$  & $1.41\times10^{4}$  & $7.07\times10^{-13}$ & $1.00$ & $3.49\times10^{-21}$ \\
$10^{-6}$  & $1.41\times10^{6}$  & $7.61\times10^{-12}$ & $1.00$ & $1.11\times10^{-16}$ \\
$10^{-8}$  & $1.41\times10^{8}$  & $4.90\times10^{-9}$  & $1.00$ & $2.22\times10^{-16}$ \\
$10^{-10}$ & $1.41\times10^{10}$ & $1.00\times10^{-10}$ & $1.00$ & $0$ \\
$10^{-12}$ & $1.41\times10^{12}$ & $1.00\times10^{-12}$ & $1.00$ & $0$ \\
$10^{-14}$ & $1.41\times10^{14}$ & $1.00\times10^{-14}$ & $1.00$ & $0$ \\
\hline
\end{tabular}
\end{table}



The orthogonal family of the adversarial proposition of the main paper isolates a single bad
pivot per $2\times2$ block at a common parameter $\varepsilon$. A variant with
geometrically decaying block parameters $\varepsilon_k=\varepsilon\,\delta^{k}$
reproduces the effect with repeated rather than isolated bad pivots, so that
the first-nonzero rule encounters a progressively worsening pivot in every
block while the matrix remains exactly orthogonal, $\kappa_2(A)=1$. On this
family the median tableau growth is $1.41\times10^{10}$ and the median scaled
defect $\|D_N\|_2/\|A\|_2$ is $1.01\times10^{-8}$ for \texttt{R0\_C0}, against
growth $1.00$ to displayed precision and defect $1.11\times10^{-16}$ for the remaining variants.
The repeated-bad-pivot construction confirms, with a different pivot geometry, the exact structural
conclusion of the single-parameter family: tableau growth cannot be bounded by $\kappa_2(A)$
alone because the matrices remain orthogonal while the first-nonzero pivots deteriorate. The large
observed inverse backward errors are a separate finite-precision statement. They force any modest
$O(u)$ backward-error constant at $\kappa_2(A)=1$ to be very large over the tested regime, but a
finite experiment does not prove the nonexistence of every possible condition-number-only
function. The empirical inequality
$\eta_{\mathrm{inv}}\le un/p^{\rm sc}_{\min}$ holds without exception on this family
(Table~\ref{tab:sm-law-by-family}).

\section{High-precision validation of the residual decomposition}
\label{sm:validation}

The decomposition $R_N=D_N-AE_{N+1}$ and the row-wise reconstruction of $D_N$ in the
main paper are algebraic consequences of the definitions and the recorded transformation factors. They are, however, statements about the algorithm as implemented,
involving the row and column permutations actually applied and the elementary
transformations actually formed, so an independent check that the instrumented
implementation realizes them is a check on the instrumentation and on the
modelling, not on the identities.

The validation bank comprised the $91\,466$ successful unrefined inverse records
of the campaign, that is, the same population from which the reliability
accounting of section~\ref{sm:reliability} is drawn. For each record an
independent verifier reproduced the working-precision pivot path (the row
order, the pivot columns, and the computed multipliers) and then evaluated,
in variable-precision arithmetic, the directly formed residual
$I-\Pi A\widehat V_{N+1}$, the defect $I-\Pi A\widetilde V_{N+1}$ obtained from
the exactly applied product of the computed transformations, the accumulated
application error $E_{N+1}$, and the row-wise reconstruction of $D_N$ from the
propagated local defects.

The tolerance requires care. The intermediate quantities include trailing
products $\widehat{\mathcal B}_{j:N}$ whose norms reach $10^{16}$ and beyond on
the Hilbert, Kahan, and Vandermonde families, so a check performed at fixed
precision, or against a fixed absolute tolerance, measures the arithmetic of the
verifier rather than the consistency of the implementation; an earlier pass in
extended-precision floating point failed for exactly this reason on those three
families, while passing to $10^{-19}$ on the families with bounded tableaux.
The reported pass therefore uses a tolerance scaled by the observed
trailing-product growth, with the working precision of the verifier set
accordingly.

Both checks passed on all $91\,219$ records for which the pivot path was
reproduced, and no algebraic discrepancy was observed. The remaining $247$
records, $0.27\%$ of the bank, were not validated: the verifier did not
reproduce the original pivot sequence, typically through a missing final pivot
in cases where the tableau is already ill-resolved in the working precision and
the selection is decided by quantities at the level of the rounding error. These
records are excluded from the validation statement and from no other analysis in
the paper.

\section{Reproducibility archive and retained bulk outputs}
\label{sec:reproducibility_archive}

The software, computational configurations, Slurm scripts, canonical campaign
manifest, compact numerical reports, provenance information, and manuscript
sources are available from
\url{https://github.com/FrancisRVillatoro/backward-error-purcell-egervary-castillo-matrix-inverse}.
The corresponding reproducibility dataset is archived on Zenodo at
\href{https://doi.org/10.5281/zenodo.21794454}
{doi:10.5281/zenodo.21794454}.

The public archive contains all $999$ recovered-task status records, an
independent certificate of the exact recovered population, compact analysis
data, the complete defect-replay record, final reports,
software-environment information, and SHA-256 provenance manifests. The
certified recovered population comprises $14{,}686{,}763$ matrices,
$146{,}867{,}630$ inverse records, and $1{,}321{,}808{,}670$ solution
records.

For each task, the recovery procedure processes the inverse and solution
streams in their deterministic generation order, retains only complete paired
matrix groups, and stops at the first incomplete or inconsistent group.
Consequently, each recovered task is a deterministic prefix of its original
task stream. The archived task-status records and the independent prefix
certificate reconstruct and verify all $999$ prefixes from the canonical
manifest, task partition, seed namespace, and matrix-identifier rule.

The complete recovered inverse and solution collection contains $999$ inverse
and $999$ solution gzip CSV shards, amounting to $1{,}998$ files and
$136{,}127{,}339{,}793$ bytes. These bulk outputs are retained by the author
for additional analyses and reviewer-requested calculations but are not
included in the public version~1.0.0 archive. Their filenames, exact byte
sizes, and SHA-256 digests are nevertheless included in the public archive.
Thus, the canonical manifest, deterministic generation rules, task-status
records, prefix certificate, and bulk-file manifests uniquely identify the
population analyzed in this work without requiring public distribution of all
bulk CSV outputs.

\bibliographystyle{siamplain}
\bibliography{castillo_refs}

\end{document}
